\section*{अध्याय \ref{ch:b1:lewis-resonance-vsepr} --- लुइस संरचनाएँ, अनुनाद और VSEPR}

\begin{solution}{exo:b1:lewis-resonance-vsepr:1}
\ce{H2O}: दो \ce{O-H} आबंध, O पर दो \omterm{def:b1:lewis-resonance-vsepr:lewis}{एकाकी युग्म}। \ce{NH3}: तीन
\ce{N-H} आबंध, N पर एक \omterm{def:b1:lewis-resonance-vsepr:lewis}{एकाकी युग्म}। \ce{CO2}: \ce{O=C=O}, हर O पर दो \omterm{def:b1:lewis-resonance-vsepr:lewis}{एकाकी
युग्म}। \ce{HCN}: \ce{H-C#N}, N पर एक \omterm{def:b1:lewis-resonance-vsepr:lewis}{एकाकी युग्म}। \ce{CH2O}: कार्बन
दो H से आबंधित और O से द्वि-आबंधित, O पर दो \omterm{def:b1:lewis-resonance-vsepr:lewis}{एकाकी युग्म}। इन पाँचों में सभी \omterm{def:b1:lewis-resonance-vsepr:formal-charge}{औपचारिक
आवेश} शून्य हैं। \ce{NH4+}: चार \ce{N-H} आबंध, कोई \omterm{def:b1:lewis-resonance-vsepr:lewis}{एकाकी
युग्म} नहीं; नाइट्रोजन $5 - 0 - 4 = +1$, जो आयन का आवेश है।
\end{solution}

\begin{solution}{exo:b1:lewis-resonance-vsepr:2}
\ce{H2S}: \ce{AX2E2}, कोणीय, $109.5^\circ$ से कम (मापा गया $92^\circ$)।
\ce{PCl3}: \ce{AX3E}, त्रिकोणीय पिरामिडी, लगभग $100^\circ$।
\ce{BF3}: \ce{AX3}, त्रिकोणीय समतलीय, $120^\circ$। \ce{SiH4}: \ce{AX4},
चतुष्फलकीय, $109.5^\circ$। \ce{CO3^{2-}}: \ce{AX3} (द्वि-आबंध
एक बार गिना जाता है), त्रिकोणीय समतलीय, $120^\circ$।
% ledger: geo:H2S.aHSH, geo:PCl3.aClPCl
\end{solution}

\begin{solution}{exo:b1:lewis-resonance-vsepr:3}
अध्रुवीय: \ce{CCl4} (चतुष्फलकीय \ce{AX4}), \ce{BF3} (\ce{AX3}),
\ce{CO2} (रैखिक \ce{AX2}): \omterm{def:b1:lewis-resonance-vsepr:dipole}{आबंध द्विध्रुव} निरस्त हो जाते हैं। ध्रुवीय: \ce{CH2Cl2}
(चतुष्फलक पर दो भिन्न प्रकार के आबंध), \ce{NH3} (पिरामिडी),
\ce{SO2} (कोणीय)।
\end{solution}

\begin{solution}{exo:b1:lewis-resonance-vsepr:4}
\ce{H+} \omterm{def:b1:lewis-resonance-vsepr:lewis-acid}{लुइस अम्ल} है और \ce{H2O} क्षारक: नया \ce{O-H} आबंध, जो
\ce{H3O+} में है, बनते समय उपसहसंयोजी है (बाद में तीनों \ce{O-H} आबंध
एक जैसे होते हैं)। \ce{Cu^{2+}} अम्ल है और हर \ce{NH3} एक क्षारक: चारों
\ce{Cu-N} आबंध उपसहसंयोजी हैं।
\end{solution}

\begin{solution}{exo:b1:lewis-resonance-vsepr:5}
\chemfig{CH_3-C(=[:60]O)-[:-60]\chemabove{O}{\scriptstyle\ominus}}
$\leftrightarrow$
\chemfig{CH_3-C(-[:60]\chemabove{O}{\scriptstyle\ominus})=[:-60]O}.
दोनों संरचनाएँ तुल्य हैं: दोनों \ce{C-O} आबंधों की कोटि 1.5 है और
लंबाई एक ही, द्वि-आबंध और एकल आबंध के बीच। मेथेनोइक
अम्ल में दोनों ऑक्सीजन भिन्न हैं (एक पर H है), संरचनाएँ
तुल्य नहीं हैं, और आबंधों की लंबाइयाँ अलग-अलग रहती हैं, 120.2 (\ce{C=O}) और
\qty{134.3}{pm} (\ce{C-OH})।
\end{solution}

\begin{solution}{exo:b1:lewis-resonance-vsepr:6}
$N = 6 + 4 \times 6 + 2 = 32$ इलेक्ट्रॉन। चार एकल आबंधों और हर ऑक्सीजन पर तीन
\omterm{def:b1:lewis-resonance-vsepr:lewis}{एकाकी युग्मों} के साथ: सल्फ़र $6 - 0 - 4 = +2$, हर ऑक्सीजन
$6 - 6 - 1 = -1$; कुल $+2 - 4 = -2$। आकृति \ce{AX4}, चतुष्फलकीय।
दो \ce{S=O} आबंधों के साथ, चार ऑक्सीजनों में से कौन-से दो
द्वि-आबंधित हों, यह चुनने पर $\binom{4}{2} = 6$ संरचनाएँ मिलती हैं।
\end{solution}

\begin{solution}{exo:b1:lewis-resonance-vsepr:7}
\ce{SF4}: \ce{AX4E}, ढेंकली। \ce{ClF3}: \ce{AX3E2}, T-आकार।
\ce{XeF2}: \ce{AX2E3}, रैखिक। \ce{XeF4}: \ce{AX4E2}, वर्ग समतलीय।
त्रिकोणीय द्विपिरामिड में अक्षीय स्थान के $90^\circ$ पर तीन पड़ोसी होते हैं,
विषुवतीय स्थान के केवल दो: भारी-भरकम \omterm{def:b1:lewis-resonance-vsepr:lewis}{एकाकी युग्म} वहाँ जाते हैं जहाँ उन्हें
सबसे कम $90^\circ$ प्रतिकर्षण झेलने पड़ें।
\end{solution}

\begin{solution}{exo:b1:lewis-resonance-vsepr:8}
सल्फ़र और फ़ॉस्फ़ोरस, ऑक्सीजन और नाइट्रोजन से बड़े और कम
विद्युत-ऋणात्मक हैं: \omterm{def:b1:lewis-resonance-vsepr:lewis}{आबंधी युग्म} केंद्रीय परमाणु से अधिक दूर रहते हैं और
हाइड्रोजन की ओर खिंचे रहते हैं, इसलिए वे एक-दूसरे को कम प्रतिकर्षित करते हैं, और \omterm{def:b1:lewis-resonance-vsepr:lewis}{एकाकी
युग्म} आबंधों को $90^\circ$ की ओर दबा देते हैं।
\end{solution}

\begin{solution}{exo:b1:lewis-resonance-vsepr:9}
\ce{CH3-C#N}: \ce{CH3} का कार्बन sp$^3$ है, नाइट्राइल कार्बन sp,
नाइट्रोजन sp; 5 $\sigma$ आबंध (तीन \ce{C-H}, \ce{C-C}, एक
\ce{C#N} में) और 2 $\pi$ आबंध। \ce{CH2=CH-CH3}: दोनों ऐल्कीन कार्बन
sp$^2$ हैं, मेथिल कार्बन sp$^3$; 8 $\sigma$ आबंध (छह \ce{C-H}, दो
\ce{C-C}) और 1 $\pi$ आबंध।
\end{solution}

\begin{solution}{exo:b1:lewis-resonance-vsepr:10}
दोनों अणु पिरामिडी हैं और नाइट्रोजन पर एक \omterm{def:b1:lewis-resonance-vsepr:lewis}{एकाकी युग्म} है। \omterm{def:b1:lewis-resonance-vsepr:lewis}{एकाकी युग्म}
एक योगदान देता है जो (ऋणात्मक से धनात्मक) \omterm{def:b1:lewis-resonance-vsepr:lewis}{एकाकी युग्म} से
नाभिक की ओर इंगित करता है। \ce{NH3} में नाइट्रोजन हर आबंध का ऋणात्मक सिरा है:
\omterm{def:b1:lewis-resonance-vsepr:dipole}{आबंध द्विध्रुव} N से हाइड्रोजनों की ओर इंगित करते हैं, \omterm{def:b1:lewis-resonance-vsepr:lewis}{एकाकी युग्म} के योगदान की
ही दिशा में, और जुड़ जाते हैं। \ce{NF3} में
फ़्लुओरीन ऋणात्मक सिरा है: \omterm{def:b1:lewis-resonance-vsepr:dipole}{आबंध द्विध्रुव} फ़्लुओरीनों से
N की ओर इंगित करते हैं, \omterm{def:b1:lewis-resonance-vsepr:lewis}{एकाकी युग्म} के योगदान के विरुद्ध, और दोनों लगभग निरस्त हो जाते हैं।
\end{solution}

\begin{solution}{exo:b1:lewis-resonance-vsepr:11}
$\mu_b = \mu/(2\cos(\theta/2)) = 1.63/(2\cos 59.75^\circ) =
\qty{1.62}{\debye} = \qty{5.40e-30}{C.m}$। तब $\delta = \mu_b/(e\,d) =
5.40 \times 10^{-30}/(1.602 \times 10^{-19} \times 143.2 \times 10^{-12})
= 0.24$।
\end{solution}

\begin{solution}{exo:b1:lewis-resonance-vsepr:12}
\ce{N=N=O}, \omterm{def:b1:lewis-resonance-vsepr:formal-charge}{औपचारिक आवेश} $-1$, $+1$, $0$ के साथ; \ce{N#N-O}, $0$,
$+1$, $-1$ के साथ। \ce{N-N} आबंध (\qty{112.8}{pm}) \ce{N2} के त्रि-आबंध
(\qty{109.8}{pm}) के पास है: संरचना \ce{N#N-O} का भार अधिक है,
और वह ऋणात्मक आवेश को ऑक्सीजन पर रखती है, जो अधिक विद्युत-ऋणात्मक
परमाणु है। फिर भी \ce{N-O} आबंध एकल आबंध से छोटा है: दूसरी
संरचना भी योगदान देती है।
\end{solution}

\begin{solution}{pb:b1:lewis-resonance-vsepr:1}
\textbf{1.} \ce{N2} 10, \ce{NO} 11, \ce{NO2} 17, \ce{N2O} 16, \ce{O3} 18,
\ce{HNO3} 24।
\textbf{2.} \ce{N#N}, हर नाइट्रोजन पर एक \omterm{def:b1:lewis-resonance-vsepr:lewis}{एकाकी युग्म} के साथ।
\textbf{3.} \ce{N=O}, ऑक्सीजन पर दो \omterm{def:b1:lewis-resonance-vsepr:lewis}{एकाकी युग्म}, नाइट्रोजन पर एक \omterm{def:b1:lewis-resonance-vsepr:lewis}{एकाकी युग्म} और एक
अकेला इलेक्ट्रॉन। 11 इलेक्ट्रॉनों, यानी विषम संख्या, के साथ किसी एक परमाणु की
गिनती सदा विषम रहती है: नाइट्रोजन के चारों ओर 7 इलेक्ट्रॉन हैं।
\textbf{4.} \ce{N=N=O}: $-1$, $+1$, $0$; \ce{N#N-O}: $0$, $+1$, $-1$।
\textbf{5.} \ce{O=O-O}: केंद्रीय ऑक्सीजन $+1$ (एक \omterm{def:b1:lewis-resonance-vsepr:lewis}{एकाकी युग्म}, तीन आबंध),
एकल-आबंधित सिरे वाला ऑक्सीजन $-1$, दूसरा 0।
\textbf{6.} \cref{ex:b1:lewis-resonance-vsepr:hno3} की तरह: नाइट्रोजन $+1$,
एकल आबंध वाला और हाइड्रोजन-रहित ऑक्सीजन $-1$, बाक़ी 0।
\textbf{7.} \ce{NO} और \ce{NO2} (11 और 17 इलेक्ट्रॉन)।
\textbf{8.} दो तुल्य संरचनाएँ, द्वि-आबंध दोनों में से किसी भी ओर:
हर आबंध की कोटि 1.5।
\textbf{9.} हाँ: \qty{127.8}{pm} द्वि-आबंध
(\qty{120.8}{pm}) और एकल आबंध (\qty{147.5}{pm}) के बीच है।
\textbf{10.} द्वि-आबंध बारी-बारी से तीनों ऑक्सीजनों में से हर एक पर; हर
\ce{N-O} आबंध तीन में से एक संरचना में द्वि-आबंध है: कोटि $4/3$।
\textbf{11.} लंबा आबंध, \qty{140.6}{pm}, \ce{N-OH} है, एक एकल आबंध।
हाइड्रोजन-रहित दोनों ऑक्सीजन दो तुल्य \omterm{def:b1:lewis-resonance-vsepr:resonance}{अनुनादी संरचनाओं} द्वारा एक $\pi$ आबंध
साझा करते हैं: कोटि 1.5 वाले दो आबंध, 121.1 और
\qty{119.9}{pm}।
\textbf{12.} नाइट्रोजन पर \omterm{def:b1:lewis-resonance-vsepr:lewis}{एकाकी युग्म} के बजाय एक अकेला इलेक्ट्रॉन है:
वह \omterm{def:b1:lewis-resonance-vsepr:lewis}{आबंधी युग्मों} को युग्म की तुलना में कम प्रतिकर्षित करता है, और आबंध
$120^\circ$ से आगे खुल जाते हैं।
\textbf{13.} \ce{N#N-O}, ऋणात्मक आवेश ऑक्सीजन पर।
\textbf{14.} \ce{O3}: \ce{AX2E}, कोणीय। \ce{N2O}: \ce{AX2} (केंद्रीय N),
रैखिक। \ce{NO3-}: \ce{AX3}, त्रिकोणीय समतलीय। \ce{HNO3} का नाइट्रोजन:
\ce{AX3}, त्रिकोणीय समतलीय।
\textbf{15.} केंद्रीय ऑक्सीजन का \omterm{def:b1:lewis-resonance-vsepr:lewis}{एकाकी युग्म} \omterm{def:b1:lewis-resonance-vsepr:lewis}{आबंधी युग्मों} को उससे अधिक
प्रतिकर्षित करता है जितना वे एक-दूसरे को करते हैं।
\textbf{16.} ओज़ोन कोणीय है और उसका केंद्रीय परमाणु सिरे वाले परमाणुओं से भिन्न है
(बीच में धनात्मक \omterm{def:b1:lewis-resonance-vsepr:formal-charge}{औपचारिक आवेश}, ऋणात्मक आवेश सिरों में बँटा हुआ):
समद्विभाजक के अनुदिश एक छोटा द्विध्रुव। \ce{N2O} रैखिक है पर उसके
दोनों सिरे भिन्न परमाणु हैं: एक छोटा द्विध्रुव। \ce{N2} सममित है: कोई नहीं।
\textbf{17.} \ce{CO2} रैखिक है (\ce{AX2}): उसके \omterm{def:b1:lewis-resonance-vsepr:dipole}{आबंध द्विध्रुव} निरस्त हो जाते हैं।
\ce{SO2} कोणीय है (\ce{AX2E}): वे जुड़ जाते हैं।
\textbf{18.} \ce{AX3E}: \omterm{def:b1:lewis-resonance-vsepr:lewis}{एकाकी युग्म} कोणों को
$109.5^\circ$ से नीचे कर देता है; मापा गया $106.7^\circ$।
\textbf{19.} \ce{AX2E2}: दो \omterm{def:b1:lewis-resonance-vsepr:lewis}{एकाकी युग्म} उसे और छोटा करते हैं; $104.5^\circ$।
\textbf{20.} $\mu = 2\mu_{OH}\cos(\theta/2)$
(\cref{prop:b1:lewis-resonance-vsepr:dipole-sum})।
\textbf{21.} $\mu_{OH} = 1.857/(2\cos 52.24^\circ) = 1.857/1.2246 =
\qty{1.52}{\debye}$।
\textbf{22.} $1.516 \times 3.336 \times 10^{-30} = \qty{5.06e-30}{C.m}$।
\textbf{23.} $e\,d = 1.602 \times 10^{-19} \times 95.8 \times 10^{-12} =
\qty{1.535e-29}{C.m}$, अर्थात् \qty{4.60}{\debye}।
\textbf{24.} $\delta = 5.06 \times 10^{-30}/1.535 \times 10^{-29} =
\textbf{0.33}$: पानी का \ce{O-H} आबंध उस आवेश का एक-तिहाई वहन करता है
जो पूर्णतः आयनिक आबंध वहन करता।
\end{solution}
